% Generated draft: requires mathematical and source review.
\documentclass[11pt]{article}
\usepackage[margin=25mm]{geometry}
\usepackage{fontspec}
\setmainfont{Latin Modern Roman}
\usepackage{amsmath,amssymb,mathtools,unicode-math}
\setmathfont{Latin Modern Math}
\usepackage{enumitem,needspace,xurl,hyperref}
\setlength{\parindent}{0pt}
\setlength{\parskip}{5pt}
\setlength{\emergencystretch}{3em}
\newcommand{\N}{\mathbb{N}}
\newcommand{\Z}{\mathbb{Z}}
\newcommand{\Q}{\mathbb{Q}}
\newcommand{\R}{\mathbb{R}}
\newcommand{\C}{\mathbb{C}}
\newcommand{\F}{\mathbb{F}}
\newcommand{\A}{\mathbb{A}}
\newcommand{\PP}{\mathbb{P}}
\newcommand{\E}{\mathbb{E}}
\newcommand{\eps}{\varepsilon}
\newcommand{\dd}{\,\mathrm{d}}
\begin{document}
\section*{UnsolvedMath ID: 701613}
\textbf{Source problem code:} catalogue entry
\subsection{Definitions and mathematical statement}
A variety of algebras of a fixed signature is the class of all algebras satisfying a specified collection of identities. A subvariety is a subclass that is itself a variety in the same signature. A semigroup is free if it is isomorphic to \(A^{+}\), the semigroup of finite nonempty words in some alphabet \(A\), with concatenation as multiplication; equivalently, each element has a unique expression as a nonempty word in the free generators.

\textbf{Supplied background.} The source states that varieties of groups form a free semigroup \(N\). It further states that varieties of linear representations over an infinite field of characteristic zero form a free semigroup \(M\), and that \(N\) acts freely on \(M\). It attributes these statements to the references listed below and says that a similar theorem holds for Lie algebras. These are reported source assertions, not independently verified results in this draft.

\textbf{Problem (B. I. Plotkin, 2018; source problem 19.65).} Find other varieties of algebras \(\Theta\), beyond the group and Lie algebra settings mentioned above, for which the same situation occurs: an analogous free semigroup of varieties, an analogous free semigroup of varieties of representations, and a free action of the former on the latter.

\textbf{Unresolved qualifications.} This is an exploratory question, and the supplied record does not specify the constructions needed to turn it into a uniquely determined formal assertion. In particular, it does not define the semigroup products, the representations appropriate to a general \(\Theta\), or the action and the meaning of its freeness. It also does not specify which varieties are included, how an identity element is treated, or precisely which parts of the group representation theorem are included in the stated Lie algebra analogue. These choices remain part of the formulation requiring review; no particular choices are imposed here.
\subsection{Short English statement}
Plotkin asks for further varieties of algebras exhibiting the free semigroup and free action phenomena reported for groups and, analogously, Lie algebras. The representation constructions, products, and meaning of a free action require clarification from the cited sources before a precise general conjecture can be stated.
\subsection{Sources}
\begin{itemize}
\item UnsolvedMath local registry, permanent ID \(701613\), exact source code \texttt{catalogue entry}, title \emph{Kourovka Notebook Problem 19.65}. The supplied registry labels the problem \emph{open}; this label is not an independently verified assessment.
\item \emph{The Kourovka Notebook (Unsolved Problems in Group Theory)}, problem 19.65, proposed by B. I. Plotkin in 2018. The supplied background identifies the 19th issue (2018) and the 21st issue edition (2026), arXiv identifier \texttt{1401.0300v48}. The supplied collection description names E. I. Khukhro and V. D. Mazurov as editors.
\item A. L. Shmel'kin, \emph{DAN SSSR}, \textbf{149} (1963), 543--545 (Russian), cited in the supplied statement for the group variety background.
\item B. H. Neuman, H. Neumann, P. M. Neumann, \emph{Math. Z.}, \textbf{80} (1962), 44--62, cited in the supplied statement for the group variety background. Author spelling is retained as supplied.
\item B. I. Plotkin, \emph{Siberian Math. J.}, \textbf{13}, no. 5 (1972), 713--729, cited in the supplied statement for the linear representation background over an infinite field of characteristic zero.
\item V. A. Parfenov, \emph{Algebra Logika}, \textbf{6}, no. 4 (1967), 61--73 (Russian), cited in the supplied statement for the Lie algebra analogue.
\item L. A. Simonyan, \emph{Siberian Math. J.}, \textbf{29}, no. 2 (1988), 276--283, cited in the supplied statement for the Lie algebra analogue.
\end{itemize}
\end{document}
